% =====================================================================
% Metadados do artigo (títulos, autoria, resumos e palavras-chave).
% Este arquivo é lido pelo LaTeX (main.tex) e pelo script build_docx.py,
% que gera as versões Word exigidas pela revista. Mantenha um comando
% por linha e o conteúdo de cada comando em uma única linha.
% =====================================================================

% --- Títulos (apenas a inicial maiúscula, salvo nomes próprios) --------
\newcommand{\tituloPT}{Fatores intracontextuais do desempenho no IDEB: segmentação socioeconômica e análise intragrupo nos municípios do Maranhão}
\newcommand{\tituloEN}{Intracontextual factors of IDEB performance: socioeconomic segmentation and within-group analysis in the municipalities of Maranhão}
\newcommand{\tituloES}{Factores intracontextuales del desempeño en el IDEB: segmentación socioeconómica y análisis intragrupo en los municipios de Maranhão}

% --- Autoria (usada apenas na versão identificada) ----------------------
% \autor{Nome completo}{Instituição de vínculo, País.}{Nota de rodapé: maior titulação. Instituição. ORCID (link). E-mail.}
% Máximo de três autores; ao menos um deve ter doutorado.
\autor{Darlon José Coqueiro}{Instituto Federal de Educação, Ciência e Tecnologia do Maranhão, Brasil.}{Graduado em Ciências da Computação pela Universidade Federal do Maranhão (UFMA). Instituto Federal de Educação, Ciência e Tecnologia do Maranhão (IFMA). ORCID: https://orcid.org/0009-0003-1296-9548. E-mail: darlon.coqueiro@ifma.edu.br.}
\autor{Orlando Donato Rocha Filho}{Instituto Federal de Educação, Ciência e Tecnologia do Maranhão, Brasil.}{Doutorado. Instituto Federal de Educação, Ciência e Tecnologia do Maranhão (IFMA). Lattes: http://lattes.cnpq.br/7455720877184126. ORCID: https://orcid.org/0000-0001-5952-7139. E-mail: orlando.rocha@ifma.edu.br.}
\autor{Danúbia Soares Pires}{Instituto Federal de Educação, Ciência e Tecnologia do Maranhão, Brasil.}{Doutorado. Instituto Federal de Educação, Ciência e Tecnologia do Maranhão (IFMA). Lattes: http://lattes.cnpq.br/4739495583287970. ORCID: https://orcid.org/0000-0002-1307-7985. E-mail: danubiapires@ifma.edu.br.}

% --- Contribuições dos autores (taxonomia CRediT) ------------------------
% \contribuicao{Nome completo (igual ao de \autor)}{Contribuições, separadas por ponto e vírgula.}
% Fonte única: aparece em informacao_autores.docx e, na versão identificada, em "Contribuição dos autores".
\contribuicao{Darlon José Coqueiro}{concepção do estudo; metodologia; coleta e curadoria dos dados; desenvolvimento dos \textit{scripts} de análise; análise dos dados; visualização; discussão dos resultados; redação do manuscrito.}
\contribuicao{Orlando Donato Rocha Filho}{orientação; revisão crítica do manuscrito; aprovação da versão final.}
\contribuicao{Danúbia Soares Pires}{orientação; revisão crítica do manuscrito; aprovação da versão final.}

% --- Resumos (100 a 250 palavras, sem siglas, fórmulas ou citações) -----
\newcommand{\resumoPT}{O desempenho educacional dos municípios maranhenses está entre os mais baixos do país, e a desigualdade socioeconômica dificulta comparações diretas entre redes de ensino. Este artigo tem como objetivo identificar os fatores escolares que diferenciam municípios de alto e de baixo desempenho no Índice de Desenvolvimento da Educação Básica de 2025, nos anos iniciais do ensino fundamental da rede municipal, entre municípios de contexto socioeconômico semelhante. A pesquisa, quantitativa e descritiva, usa o Censo Escolar de 2025 e o Censo Demográfico de 2022 para os 216 municípios do Maranhão com o índice divulgado. O algoritmo \textit{K-Means} agrupa os municípios em três grupos, pela renda \textit{per capita} e pela proporção de moradores em pobreza extrema; nos grupos com amostra suficiente, o teste de Mann-Whitney, com tamanho de efeito e correção para comparações múltiplas, compara os quartis superior e inferior de desempenho. Nos grupos de renda baixa e intermediária, a proporção de docentes com ensino superior e com licenciatura está associada ao desempenho, com efeito médio; no intermediário, destaca-se também a quadra esportiva. Conclui-se que a formação docente é o fator escolar mais consistentemente associado ao desempenho, embora o desenho associativo não permita estimar seu efeito causal.}
\newcommand{\palavrasPT}{desempenho educacional; formação docente; eficácia escolar.}

\newcommand{\resumoEN}{The educational performance of municipalities in Maranhão is among the lowest in Brazil, and socioeconomic inequality hinders direct comparisons between school systems. This article aims to identify the school factors that distinguish high- and low-performing municipalities in the 2025 Basic Education Development Index, in the early years of elementary education in municipal schools, among municipalities with similar socioeconomic contexts. The study, quantitative and descriptive, uses the 2025 School Census and the 2022 Demographic Census for the 216 municipalities of Maranhão for which the index was released. The K-Means algorithm groups municipalities into three groups by per capita income and the share of residents in extreme poverty; within the groups with sufficient sample size, the Mann-Whitney test, with effect sizes and correction for multiple comparisons, compares the upper and lower performance quartiles. In the low- and middle-income groups, the share of teachers with higher education and with a teaching degree is associated with performance, with a medium effect size; in the middle-income group, sports courts also stand out. It is concluded that teacher education is the school factor most consistently associated with performance, although the associative design does not allow its causal effect to be estimated.}
\newcommand{\palavrasEN}{educational performance; teacher education; school effectiveness.}

\newcommand{\resumoES}{El desempeño educativo de los municipios de Maranhão está entre los más bajos de Brasil, y la desigualdad socioeconómica dificulta las comparaciones directas entre redes de enseñanza. Este artículo tiene como objetivo identificar los factores escolares que diferencian a los municipios de alto y de bajo desempeño en el Índice de Desarrollo de la Educación Básica de 2025, en los años iniciales de la enseñanza fundamental de la red municipal, entre municipios de contexto socioeconómico similar. La investigación, cuantitativa y descriptiva, usa el Censo Escolar de 2025 y el Censo Demográfico de 2022 para los 216 municipios de Maranhão con el índice divulgado. El algoritmo \textit{K-Means} agrupa los municipios en tres grupos, por la renta \textit{per cápita} y la proporción de residentes en pobreza extrema; en los grupos con muestra suficiente, la prueba de Mann-Whitney, con tamaño del efecto y corrección por comparaciones múltiples, compara los cuartiles superior e inferior de desempeño. En los grupos de renta baja e intermedia, la proporción de docentes con educación superior y con licenciatura se asocia al desempeño, con efecto medio; en el grupo intermedio, destacan también las canchas deportivas. Se concluye que la formación docente es el factor escolar asociado de forma más consistente al desempeño, aunque el diseño asociativo no permite estimar su efecto causal.}
\newcommand{\palavrasES}{desempeño educativo; formación docente; eficacia escolar.}
